\capredactadotrue
\section{Segmentación}

La segmentación del mercado de \nombreproyecto{} tiene como propósito identificar y agrupar a los usuarios según sus necesidades, su rol dentro de la dinámica de juego y su comportamiento de uso. Al tratarse de una herramienta tecnológica orientada a un nicho específico (los juegos de rol de mesa), definir correctamente estos segmentos es fundamental para orientar las prioridades de desarrollo del producto mínimo viable (MVP) y focalizar los esfuerzos de captación inicial.

\subsection{Segmentación de consumidores y negocios}
\nombreproyecto{} opera bajo un modelo centrado en el consumidor final (B2C) para las funcionalidades base de la plataforma de software (SaaS), complementado por un esquema colaborativo de usuario a usuario (C2C) a través del \textit{marketplace} integrado. Por lo tanto, la segmentación se aborda exclusivamente desde la óptica de consumidores individuales, y se prescinde de una segmentación corporativa (B2B).

Para definir el mercado meta se aplican las siguientes variables sobre el consumidor final.

\subsubsection{Segmentación geográfica}
En la etapa inicial de lanzamiento y validación, el segmento se concentra en usuarios residentes en Argentina, con especial foco en los polos urbanos de mayor actividad de comunidades de rol: principalmente Rosario, Córdoba y Buenos Aires. Por la naturaleza digital del producto, el segmento tiene proyección de escalabilidad hacia el resto de América Latina y el mercado hispanohablante en general.

\subsubsection{Segmentación demográfica}
El público objetivo principal está formado por adolescentes y adultos jóvenes de entre 16 y 35 años. Este grupo es el núcleo demográfico con mayor ingreso reciente al pasatiempo, impulsado por la exposición de los juegos de rol en medios digitales, en videojuegos (como \textit{Baldur's Gate 3}) y en la cultura popular contemporánea.

\subsubsection{Segmentación psicográfica}
Se dirige a personas con afinidad por las actividades recreativas grupales de carácter presencial que, al mismo tiempo, poseen un alto nivel de alfabetización digital y de adopción tecnológica. Son usuarios que valoran la creatividad, la inmersión narrativa y la cooperación, y que están muy familiarizados con las comunidades en línea, las redes sociales y las plataformas de comunicación instantánea (como Discord o Twitch).

\subsubsection{Segmentación conductual}
Comprende a quienes participan en juegos de rol de mesa, ya sea de forma frecuente (campañas regulares) o esporádica. Ante la necesidad que el producto resuelve, este segmento experimenta una alta frustración por la fragmentación de la información (manuales físicos combinados con aplicaciones de notas aisladas o grupos de mensajería), que enlentece sus partidas. Busca herramientas que reduzcan la carga administrativa y simplifiquen la organización sin restar valor a la experiencia de estar frente a una mesa física con sus amigos.

\subsection{Grupos diferenciados de consumidores}
Dado que el modelo es B2C con un componente C2C, los grupos se identifican desde la lógica de consumidores individuales y no de empresas. No obstante, los creadores de contenido del \textit{marketplace} pueden analizarse con una lógica híbrida, ya que actúan como pequeños oferentes dentro de la plataforma.

Los grupos se diferencian principalmente por:
\begin{itemize}
  \item \textbf{Rol dentro del juego:} \textit{Dungeon Master}, jugador o creador de contenido.
  \item \textbf{Nivel de experiencia:} jugadores nuevos (ingresantes al pasatiempo) frente a jugadores experimentados.
  \item \textbf{Frecuencia de juego:} grupos con sesiones regulares (semanales o quincenales) frente a grupos esporádicos.
  \item \textbf{Ubicación geográfica:} mercado inicial (Argentina) frente a mercados de expansión futura (América Latina e internacional).
\end{itemize}

A partir de estos criterios se identifican cuatro grupos:
\begin{itemize}
  \item \textbf{DMs experimentados:} gestionan varias campañas de forma simultánea o recurrente, tienen una alta necesidad de organización y mayor propensión a pagar por herramientas que reduzcan la carga administrativa.
  \item \textbf{DMs novatos o ingresantes al pasatiempo:} enfrentan la elevada curva de aprendizaje de las herramientas del sector (véase \ref{sec:fortalezas-debilidades-sector}); valoran especialmente la simplicidad de uso y el acompañamiento durante la planificación de sus primeras campañas.
  \item \textbf{Jugadores activos:} participan de campañas ya creadas y requieren menos funcionalidades avanzadas, pero les interesa el acceso rápido a la información de sus personajes.
  \item \textbf{Creadores de contenido:} usuarios con intereses narrativos o creativos (construcción de mundos, diseño de mapas) que buscan un canal de distribución y una eventual monetización de su contenido.
\end{itemize}

Todos los grupos comparten el acceso a dispositivos móviles o computadoras y la familiaridad con plataformas digitales y comunidades en línea, en coherencia con el perfil habitual de la comunidad TTRPG.

\subsection{Potenciales usuarios y compradores}
Se identifican tres perfiles de usuario, cada uno con un rol distinto dentro del modelo de negocio:
\begin{itemize}
  \item \textbf{\textit{Dungeon Masters} (DM):} usuario principal y de mayor involucramiento, responsable de administrar campañas, sesiones y contenido. Es quien enfrenta la mayor carga administrativa y, por lo tanto, el principal beneficiario de la propuesta de valor.
  \item \textbf{Jugadores:} usuarios secundarios que participan de las campañas creadas por un DM, consultan la información de sus personajes y colaboran en la organización de las partidas.
  \item \textbf{Creadores de contenido:} usuarios que ofrecen campañas, mapas y recursos narrativos en el \textit{marketplace}; constituyen la oferta de contenido y, a la vez, una fuente adicional de ingresos para la plataforma.
\end{itemize}

En términos de compra, el comprador principal, tanto por la conversión a planes premium como por el consumo dentro del \textit{marketplace}, es el DM, pues administra la campaña y toma la mayoría de las decisiones sobre qué herramientas y contenido incorporar. Los jugadores, aunque son usuarios activos, tienen un rol de comprador secundario, vinculado sobre todo a la personalización de sus propios personajes.

Los creadores de contenido son un tercer perfil de comprador y vendedor dentro de la lógica C2C: no utilizan la plataforma como un jugador o un DM, sino como canal de distribución y monetización de su propio contenido.

\subsection{Pautas de comportamiento esperado por segmento}
\begin{itemize}
  \item \textbf{DMs:} uso recurrente y de alta frecuencia, generalmente asociado a la cadencia de las sesiones de su grupo (semanal o quincenal). Lo motiva la necesidad de reducir las tareas administrativas y centralizar la información. Tienen la mayor propensión a explorar y, eventualmente, pagar por funcionalidades premium.
  \item \textbf{Jugadores:} uso más esporádico, acotado a la consulta de su personaje y a eventos puntuales de la sesión. Los motivan la comodidad y el acceso rápido a la información, con baja fricción de uso.
  \item \textbf{Creadores de contenido:} uso orientado a la publicación y actualización periódica de recursos en el \textit{marketplace}. Su motivación es mixta: reconocimiento dentro de la comunidad y monetización del contenido creado.
\end{itemize}
